% ============================================================
%  PREAMBOLO DELLA DISPENSA
%  Pacchetti extra, colori dei livelli e componenti TikZ riusabili
% ============================================================

\usepackage{booktabs}
\usepackage{multirow}
\usepackage{array}
\usepackage{makecell}
\usepackage[font=small,labelfont={bf,color=primary!85!black},labelsep=endash]{caption}
\usepackage{pgfplots}
\pgfplotsset{compat=1.18}

\usetikzlibrary{shapes.symbols, shapes.misc, shapes.arrows, shadows, fit,
  backgrounds, decorations.pathmorphing, decorations.markings,
  decorations.pathreplacing, patterns, chains, matrix, arrows.meta, calc,
  positioning}

\renewcommand{\arraystretch}{1.25}

% Titoli dei riquadri protetti da graffe: così virgole e "=" nel titolo non rompono tcolorbox
\renewcommand{\info}[2]{\begin{infobox}{{#1}}#2\end{infobox}}
\renewcommand{\warning}[2]{\begin{warningbox}{{#1}}#2\end{warningbox}}
\renewcommand{\error}[2]{\begin{errorbox}{{#1}}#2\end{errorbox}}

% Nei listati il simbolo $ (prompt della shell) non deve attivare la matematica
\gdef\documentlstlabel{Listato}
\gdef\documentlsttitle{Listati}
\lstset{mathescape=false, upquote=true, title={}}
\captionsetup[lstlisting]{labelformat=empty, font={footnotesize,sf}, textfont={color=primary!80!black}, skip=2pt}
\lstdefinelanguage{http}{morekeywords={GET,POST,HEAD,PUT,DELETE,HTTP/1.1,HTTP/1.0}, sensitive=true}

% ------------------------------------------------------------
% Colori associati ai livelli della pila (usati in tutte le figure)
% ------------------------------------------------------------
\definecolor{lapp}{HTML}{F4B183}   % applicazione
\definecolor{lpres}{HTML}{F8CBAD}  % presentazione
\definecolor{lsess}{HTML}{FBE5D6}  % sessione
\definecolor{ltra}{HTML}{A9D18E}   % trasporto
\definecolor{lnet}{HTML}{9DC3E6}   % rete
\definecolor{llink}{HTML}{FFE699}  % collegamento
\definecolor{lphy}{HTML}{D9D9D9}   % fisico
\definecolor{netblue}{HTML}{2B50AA}
\definecolor{netred}{HTML}{C00000}
\definecolor{netgreen}{HTML}{2E7D32}
\definecolor{netorange}{HTML}{E46C0A}
\definecolor{cloudbg}{HTML}{DDEBF7}

% ------------------------------------------------------------
% Icone dei dispositivi (simboli in stile Cisco)
%   \icon[larghezza]{router|switch|hub|accesspoint|pc|laptop|server|smartphone}
% ------------------------------------------------------------
\newcommand{\icon}[2][0.9cm]{\includegraphics[width=#1]{icons/#2}}

% ------------------------------------------------------------
% Stili TikZ
% ------------------------------------------------------------
\tikzset{
  % scatola generica
  box/.style={draw=black!70, rounded corners=2pt, minimum height=0.8cm,
    minimum width=2.2cm, align=center, fill=white, font=\small},
  % strato di una pila protocollare
  layer/.style={draw=black!60, minimum height=0.62cm, minimum width=2.8cm,
    align=center, font=\small\sffamily},
  lay app/.style={layer, fill=lapp},
  lay pres/.style={layer, fill=lpres},
  lay sess/.style={layer, fill=lsess},
  lay tra/.style={layer, fill=ltra},
  lay net/.style={layer, fill=lnet},
  lay link/.style={layer, fill=llink},
  lay phy/.style={layer, fill=lphy},
  % nodo di un grafo
  gnode/.style={circle, draw=netblue!80!black, fill=netblue!60, minimum size=0.42cm,
    inner sep=0pt},
  hostnode/.style={circle, draw=netgreen!70!black, fill=netgreen!45, minimum size=0.42cm,
    inner sep=0pt},
  % collegamenti
  wired/.style={line width=1.1pt, netblue},
  wireless/.style={line width=1.1pt, netblue, dashed, dash pattern=on 2pt off 2pt},
  backbone/.style={line width=3pt, netblue!85!black},
  % nuvola (rete / Internet)
  netcloud/.style={cloud, cloud puffs=11, cloud puff arc=110, aspect=2,
    draw=netblue!60, fill=cloudbg, align=center, font=\small\sffamily},
  % freccia di un messaggio nei diagrammi temporali
  msg/.style={-{Stealth[length=2.2mm]}, thick},
  msglost/.style={-{Rays[n=6]}, thick, netred},
  % etichette piccole
  lbl/.style={font=\footnotesize\sffamily, align=center},
  smalllbl/.style={font=\scriptsize\sffamily, align=center},
  % campi di un header (pacchetti, segmenti, frame)
  field/.style={draw=black!70, minimum height=0.65cm, align=center, font=\scriptsize\sffamily,
    inner sep=1pt},
  % evidenziazione di un percorso
  pathhl/.style={line width=2.2pt, netred},
}

% Linea temporale verticale per diagrammi di scambio messaggi:
%   \timeline{nome}{x}{etichetta}{lunghezza}
\newcommand{\timeline}[4]{%
  \node[lbl, font=\small\bfseries\sffamily] (#1) at (#2,0.35) {#3};
  \draw[thick, black!70] (#2,0) -- (#2,-#4);
}

% Riquadro "In breve" per i riepiloghi a fine capitolo
\newcommand{\riepilogo}[1]{%
  \gbox{\iinfo\ In breve}{primary!80!black}{#1}}

% Numero di passo evidenziato (usabile anche dentro i nodi TikZ)
\newcommand{\circled}[1]{{\setlength{\fboxsep}{1.2pt}\colorbox{primary}{\textcolor{white}{\bfseries\scriptsize\sffamily\,#1\,}}}}

% Evidenziazione di un termine chiave in linea
\newcommand{\key}[1]{\textbf{\color{primary!85!black}#1}}

% Intestazione dello svolgimento negli esercizi della Parte VII
\newcommand{\svolgimento}{\par\smallskip\noindent{\sffamily\bfseries\color{primary!80!black}Svolgimento}\par\nobreak\smallskip}

% Esempi ed esercizi numerati per capitolo (le sezioni non sono numerate)
\counterwithin*{esempio}{chapter}
\counterwithin*{esercizio}{chapter}
% Esempi ed esercizi composti in una scatola: se stanno in una pagina non si
% spezzano mai (passano interi alla pagina successiva); se sono più lunghi di
% una pagina la scatola viene "svuotata" e si spezzano normalmente.
\newbox\rpdesbox
\newcommand{\rpdemetti}{%
  \ifdim\dimexpr\ht\rpdesbox+\dp\rpdesbox\relax>\textheight
    \unvbox\rpdesbox
  \else
    \box\rpdesbox
  \fi}
\renewenvironment{esempio}[1]{%
  \par\addvspace{0.8em}%
  \setbox\rpdesbox\vbox\bgroup
  \refstepcounter{esempio}%
  \noindent\rule{\linewidth}{0.5pt}\par%
  \vspace{0.15em}%
  \noindent\textbf{Esempio \arabic{esempio}}%
  \ifx\relax#1\relax\else\textbf{ --- #1}\fi\par%
  \vspace{0.3em}%
}{%
  \par\vspace{0.3em}%
  \noindent\rule{\linewidth}{0.5pt}\par
  \egroup
  \rpdemetti
  \par\vspace{0.8em}%
}
\renewenvironment{esercizio}[1]{%
  \par\addvspace{0.8em}%
  \setbox\rpdesbox\vbox\bgroup
  \refstepcounter{esercizio}%
  \noindent\rule{\linewidth}{1.2pt}\vspace{1pt}\par%
  \noindent\rule{\linewidth}{0.4pt}\par%
  \vspace{0.15em}%
  \noindent\textbf{Esercizio \arabic{esercizio}}%
  \ifx\relax#1\relax\else\textbf{ --- #1}\fi\par%
  \vspace{0.3em}%
}{%
  \par\vspace{0.3em}%
  \noindent\rule{\linewidth}{0.4pt}\vspace{1pt}\par%
  \noindent\rule{\linewidth}{1.2pt}\par
  \egroup
  \rpdemetti
  \par\vspace{0.8em}%
}

% Niente stiramento verticale delle pagine (evita buchi prima di blocchi lunghi)
\raggedbottom

% Pagine di apertura delle parti nello stile del documento
\titleformat{\part}[display]
  {\centering\color{primary}}
  {\fontsize{20}{24}\selectfont\Gobold Parte \thepart}
  {18pt}
  {\color{primary!80!black}\fontsize{32}{40}\selectfont\Gobold #1}
  [\vspace{10pt}{\color{primary!40}\rule{0.35\linewidth}{2pt}}]
\assignpagestyle{\part}{empty}

% Omini stilizzati per i diagrammi (ispirati ai disegni alla lavagna)
\tikzset{
  persona/.pic={
    \draw[thick] (0,0.58) circle (0.13);
    \draw[thick] (0,0.45) -- (0,0.05);
    \draw[thick] (-0.22,0.3) -- (0,0.36) -- (0.22,0.3);
    \draw[thick] (0,0.05) -- (-0.16,-0.32);
    \draw[thick] (0,0.05) -- (0.16,-0.32);
  },
  personacapelli/.pic={
    \pic{persona};
    \draw[thick] (-0.12,0.66) .. controls (-0.3,0.5) and (-0.22,0.3) .. (-0.26,0.2);
    \draw[thick] (0.12,0.66) .. controls (0.3,0.5) and (0.22,0.3) .. (0.26,0.2);
  },
  pinint/.style={circle, draw, inner sep=0pt, minimum size=4.5pt, fill=white},
}
